\subsection{রোবটিকস}

\begin{learningbox}
এই অংশ শেষে শিক্ষার্থী পারবে:
\begin{itemize}
\item রোবট ও রোবটিকসের ধারণা ব্যাখ্যা করতে।
\item রোবটের sensor, actuator, controller ও power system-এর কাজ লিখতে।
\item শিল্প, চিকিৎসা, শিক্ষা, কৃষি ও প্রতিরক্ষায় রোবটের ব্যবহার চিহ্নিত করতে।
\item রোবট ব্যবহারের সুবিধা ও সীমাবদ্ধতা বিশ্লেষণ করতে।
\end{itemize}
\end{learningbox}

\subsubsection{রোবটিকসের ধারণা}

রোবটিকস হলো রোবটের নকশা, তৈরি, নিয়ন্ত্রণ, programming এবং ব্যবহার নিয়ে কাজ করা প্রযুক্তিক্ষেত্র। রোবট সাধারণত sensor দিয়ে পরিবেশ থেকে তথ্য নেয়, controller দিয়ে সিদ্ধান্ত নেয় এবং actuator বা motor দিয়ে কাজ সম্পন্ন করে।

\begin{definitionbox}{সংজ্ঞা}
রোবটিকস হলো engineering ও computer science-এর সমন্বিত শাখা, যেখানে programmable machine বা robot তৈরি ও নিয়ন্ত্রণ করা হয়, যাতে তা স্বয়ংক্রিয় বা আধা-স্বয়ংক্রিয়ভাবে নির্দিষ্ট কাজ করতে পারে।
\end{definitionbox}

\begin{mistakebox}{সাধারণ ভুল}
সব machine রোবট নয়। কোনো machine যদি শুধু switch চাপলে একই কাজ করে, তবে সেটি automation হতে পারে। কিন্তু রোবট সাধারণত programmable, sensor-based এবং নির্দিষ্ট মাত্রায় পরিবেশের সঙ্গে interact করতে পারে।
\end{mistakebox}

\subsubsection{রোবটের প্রধান অংশ}

\begin{center}
\begin{tabular}{|p{0.28\textwidth}|p{0.58\textwidth}|}
\hline
\textbf{অংশ} & \textbf{কাজ} \\
\hline
Sensor & আলো, তাপ, দূরত্ব, শব্দ, চাপ, ছবি বা অবস্থান শনাক্ত করে \\
\hline
Controller & sensor data বিশ্লেষণ করে command দেয় \\
\hline
Actuator & motor, arm, wheel বা gripper নড়াচড়া করায় \\
\hline
Power system & battery বা electric supply দিয়ে শক্তি দেয় \\
\hline
End effector & robot arm-এর শেষ অংশ; ধরা, কাটা, welding বা painting করে \\
\hline
Program & robot কীভাবে কাজ করবে তা নির্ধারণ করে \\
\hline
\end{tabular}
\end{center}

\subsubsection{রোবটের বৈশিষ্ট্য}

\begin{keypointbox}{প্রধান বৈশিষ্ট্য}
\begin{itemize}
\item নির্দিষ্ট কাজ বারবার একইভাবে করতে পারে।
\item বিপজ্জনক পরিবেশে মানুষের বদলে কাজ করতে পারে।
\item sensor data নিয়ে পরিবেশ অনুযায়ী response দিতে পারে।
\item programming বদলালে কাজের ধরন বদলানো যায়।
\item AI যুক্ত হলে কিছু ক্ষেত্রে সিদ্ধান্ত নেওয়ার ক্ষমতা বাড়ে।
\end{itemize}
\end{keypointbox}

\subsubsection{রোবটের ব্যবহারক্ষেত্র}

\begin{examplebox}{শিল্পকারখানা}
গাড়ি assembly line-এ robot welding, painting, lifting ও quality inspection করতে পারে। এতে উৎপাদন দ্রুত হয় এবং মানুষের জন্য ঝুঁকিপূর্ণ কাজ কমে।
\end{examplebox}

\begin{examplebox}{চিকিৎসা}
Surgical robot ডাক্তারকে সূক্ষ্ম operation-এ সহায়তা করতে পারে। Rehabilitation robot রোগীর হাত-পা নড়াচড়া practice করাতে পারে।
\end{examplebox}

\begin{keypointbox}{আরও ব্যবহার}
\begin{itemize}
\item কৃষি: crop monitoring, spraying, harvesting।
\item শিক্ষা: STEM lab, coding practice, robotics competition।
\item প্রতিরক্ষা: bomb disposal, surveillance, drone operation।
\item মহাকাশ: rover, robotic arm, satellite maintenance।
\item গৃহস্থালি: robotic vacuum cleaner, assistive robot।
\end{itemize}
\end{keypointbox}

\subsubsection{রোবটিকসের সুবিধা}

\begin{itemize}
\item ঝুঁকিপূর্ণ কাজ মানুষের বদলে করতে পারে।
\item একই কাজ দ্রুত ও নির্ভুলভাবে বারবার করতে পারে।
\item উৎপাদনশীলতা ও quality control বাড়ায়।
\item ২৪ ঘণ্টা কাজ করতে পারে, fatigue কম।
\item চিকিৎসা ও গবেষণায় সূক্ষ্ম কাজের সহায়তা করে।
\end{itemize}

\subsubsection{রোবটিকসের সীমাবদ্ধতা}

\begin{itemize}
\item installation ও maintenance cost বেশি হতে পারে।
\item skilled programmer ও technician দরকার।
\item ভুল program বা sensor error হলে দুর্ঘটনা ঘটতে পারে।
\item কিছু repetitive job কমে যেতে পারে।
\item মানবিক বিচার, empathy ও সামাজিক context রোবট সহজে বুঝতে পারে না।
\end{itemize}

\begin{boardbox}{বোর্ড ফোকাস}
রোবটিকসের উত্তরে sensor, controller, actuator—এই তিনটি অংশ উল্লেখ করলে উত্তর বেশি নির্ভুল হয়। ব্যবহারক্ষেত্র লিখতে শিল্প, চিকিৎসা, প্রতিরক্ষা ও মহাকাশ থেকে উদাহরণ দাও।
\end{boardbox}

\begin{practicebox}
\textbf{MCQ}
\begin{enumerate}
\item রোবটের পরিবেশ শনাক্তকারী অংশ কোনটি?\\
ক. Sensor \quad খ. Monitor \quad গ. Printer \quad ঘ. Scanner only
\item Robot arm নড়াচড়া করায় কোন অংশ?\\
ক. Actuator \quad খ. ROM \quad গ. Browser \quad ঘ. Router
\end{enumerate}

\textbf{সংক্ষিপ্ত প্রশ্ন}
\begin{enumerate}
\item রোবটিকস বলতে কী বোঝ?
\item রোবট ব্যবহারের দুটি সুবিধা লেখ।
\end{enumerate}
\end{practicebox}

